\section{Introduction}
Every year, American communities hold thousands of elections on one question: may the local government borrow. These bond referenda are, by simple count, among the most numerous policy elections in the United States, and the rules that govern them are old. California has required a supermajority for local general obligation debt since 1879. Kentucky wrote its consent requirement into Section 157 of the 1891 constitution. The machinery is so familiar that it is rarely studied as machinery. This paper asks what it does: how the rules that make voters consent to debt shape what communities ask, what they build, and who pays.

Start with what such a rule is. A requirement that debt be approved by a majority, or by fifty-five per cent, or by two thirds of those voting, is not just a hurdle. It defines a public: whose agreement counts, and how much of it must be gathered, before a community may pledge its taxing power. That definition has never been neutral. For much of the twentieth century, several states limited the bond franchise to property owners in constitutional text; the Supreme Court struck the freeholder qualification around 1970 [V]. Its demography did not fall with it. Bond elections are commonly held at times when those who turn out are older and more propertied than the community they decide for, so the public the law once defined survives in practice (Anzia 2014 [V]; Hajnal, Kogan and Markarian 2022 [V]). And beside this consenting public stands a second public the rules create by leaving it out: the ratepayers, tenants and users who repay the forms of debt that need no one's consent, and who hold a veto over none of it. The same Texas majority rule that made Harris County assemble 511,375 yes votes lets a developer's utility district borrow on two.

The paper's argument is short. An authorisation rule makes two choices. It sets scope: which debt must be voted on at all. And it sets height: how much of which electorate must agree. Governments respond at two points. Before any vote, they choose the channel, and they shape what reaches the ballot: how much to ask, how many purposes to bundle, when to schedule. After a refusal, they wait and ask again, or they route the money through forms that need no public. Because only goods that can bill their users can take that second route, the rules do not fall evenly. Water can leave the ballot; schools cannot. What escapes the voters accumulates out of their sight, financed by flat fees rather than taxes.

Three bodies of evidence carry the argument. First, the authorisation system is measured, for the first time, from the documents it produces. Nearly every bond offering document states the authority behind the issue, so the voted share of borrowing can be observed directly, in all fifty states, rather than guessed from statutes: 93.7 per cent of documents yield a determination, and where the documents cite elections, the citations match independent election records almost exactly where the records are complete. Second, the causes are identified at the thresholds themselves. Across 11,889 bond measures in five states, a measure that barely clears its statutory bar behaves very differently from one that barely misses: 24 points more likely to issue in the vote year itself, twice the borrowing, and, in Census spending data, construction falling nearly one-for-one with the refused debt. Third, every refusal is followed forward: who returns, who reroutes, who gives up.

What the evidence shows can be said in three sentences. The vote is real: authorisation moves borrowing and building, not paperwork. It is a queue and a sorter, not a wall: refused governments return at the next election with the same request and usually win, fewer than one in five blocked projects die, and the requirement's main work is deciding when things are built, through which instrument, and at whose expense. And it consults a particular public: the effect is three times larger where homeowners dominate, absent where the community is diverse, unrelated to income or party, and the majorities blocked by supermajority bars sit in poorer school districts. The rules of 1879 still answer to the electorate of 1879.

The paper's relations to existing research are developed inside the argument (Section 2); stated briefly: studies of direct democracy have concentrated on the votes themselves, and the findings here locate the institution's work before any vote and after any refusal; research on fiscal limits has treated substitution as leakage, and here it is the institution's equilibrium and the reason it survives; and the study of American political development gains a persistence mechanism, selective escape, and a measured answer to who the oldest consent institution in American fiscal law still serves.

